\documentclass[11pt]{article}

\usepackage[T1]{fontenc}
\usepackage{lmodern}
\usepackage[margin=1in]{geometry}
\usepackage{amsmath,amssymb,amsthm}
\usepackage{booktabs}
\usepackage[hidelinks]{hyperref}

\newtheorem{theorem}{Theorem}
\newtheorem{lemma}[theorem]{Lemma}
\theoremstyle{remark}
\newtheorem{remark}[theorem]{Remark}

\newcommand{\F}{\mathbb{F}}
% The bound proved. 34 needs both solvers to exclude M = 33; 33 needs M = 32 only.
\newcommand{\newbound}{34}
\newcommand{\lastM}{33}

\title{A new lower bound for the mixed covering code $K(2,5,1)$}
\author{Noah Parsons\\ \small Independent researcher\thanks{%
\texttt{parsons.m.noah@gmail.com}, ORCID \href{https://orcid.org/0009-0000-7224-6040}{0009-0000-7224-6040}.}}
\date{October 2026}

\begin{document}
\maketitle

\begin{abstract}
Let $K(t,b,1)$ be the minimum size of a code in $\F_3^t\times\F_2^b$ with covering radius $1$.
The tables of mixed binary/ternary covering codes give $32\le K(2,5,1)\le 36$, the smallest open
entry for covering radius $1$. We prove $K(2,5,1)\ge\newbound$. The space is split into nine
blocks by the two ternary coordinates; a counting inequality leaves few possible distributions of
codewords over the blocks, and each remaining case is shown infeasible by integer programming.
Two independent solvers agree on every case, and every step of the reduction is checked by
separately written code.
\end{abstract}

\section{Introduction}

A code $C\subseteq\F_3^t\times\F_2^b$ has \emph{covering radius $1$} if every word is within
Hamming distance $1$ of some codeword, and $K(t,b,1)$ is the minimum size of such a
code~\cite{cohen-book}. Bounds for $t+b\le 13$ are collected in the tables of
K\'eri~\cite{keri}, following~\cite{oh,bow}. For $K(2,5,1)$, the entry with the smallest space
($9\cdot 32=288$ words) among the open radius-$1$ entries, the tables give
\[
32\le K(2,5,1)\le 36 .
\]
The upper bound is due to Hämäläinen and Rankinen~\cite{hr}, the lower bound to Lang, Quistorff
and Schneider~\cite{lqs}. We are not aware of any later improvement.

\begin{theorem}\label{thm:main}
$K(2,5,1)\ge\newbound$.
\end{theorem}

Since a code with fewer than $M$ words can be enlarged to exactly $M$ words without losing the
covering property, it suffices to show that no covering code with exactly $M$ words exists for
$M=32$ and $M=\lastM$.

\section{Method}

Write words as $(a,b,r)$ with $a,b\in\F_3$ and $r\in\F_2^5$. The \emph{block} $c=(a,b)$ is the
set of $32$ words with these ternary coordinates; the nine blocks form a $3\times 3$ grid. For a
code $C$ let $M_c=|C\cap c|$.

\begin{lemma}\label{lem:count}
If $C$ has covering radius $1$, then for every block $c$,
\[
6M_c+\sum_{d\in N(c)} M_d\ \ge\ 32,
\]
where $N(c)$ is the set of the four blocks in the same row or column of the grid as $c$.
\end{lemma}

\begin{proof}
Every word of $c$ lies in the ball of some codeword, so the balls of radius $1$ around the
codewords, intersected with $c$, cover all $32$ words of $c$. A codeword in $c$ meets $c$ in
itself and its $5$ binary neighbours. A codeword in a block of $N(c)$ differs from $c$ in exactly
one ternary coordinate, so its ball meets $c$ in exactly one word. A codeword in any other block
differs from every word of $c$ in both ternary coordinates, so its ball misses $c$. Summing gives
the inequality.
\end{proof}

The search for a code with $M$ words therefore proceeds in two steps.
\begin{enumerate}
\item \textbf{Distributions.} List every $(M_c)_c$ with $0\le M_c\le 32$, $\sum_c M_c=M$, that
  satisfies Lemma~\ref{lem:count} for all nine blocks. Permuting the symbols of either ternary
  coordinate, and swapping the two ternary coordinates, are isometries; they act on the grid as a
  group of order $72$, and one representative of each orbit of distributions is kept.
\item \textbf{Exact cases.} For each representative, decide by integer programming whether there
  is a covering code with exactly $M_c$ codewords in each block $c$: one binary variable per word,
  one covering constraint per word, and one count constraint per block.
\end{enumerate}

\section{Computation}

\begin{table}[h]
\centering
\begin{tabular}{rrrccc}
\toprule
$M$ & distributions & up to symmetry & SCIP 10.0 & Gurobi 13.0 & HiGHS \\
\midrule
32 & 306   & 10 & 10 infeasible & 10 infeasible & 10 infeasible \\
33 & 1\,194 & 32 & 32 infeasible & 32 infeasible & 31 infeasible, 1 open \\
\bottomrule
\end{tabular}
\caption{Cases for $K(2,5,1)$. Each solver was run on a model built by separate code. HiGHS
(through SciPy~1.16.3) is a third, slower check; one case at $M=33$ had not finished after
30 minutes.}
\label{tab:cases}
\end{table}

Table~\ref{tab:cases} summarises the computation. Every case was found infeasible by SCIP~10.0
(through PySCIPOpt) and, with a separately written model, by Gurobi~13.0 (through gurobipy),
single-threaded on one core of an Intel Core Ultra~7 255HX. Each case took at most $33$ seconds
with SCIP and at most $1.5$ seconds with Gurobi. HiGHS, a third solver, agrees on every case it
finished.

The following checks were made on the steps before the solvers.
\begin{itemize}
\item \emph{Lemma~\ref{lem:count} from the definition.} For each block $c$ and each word $u$, the
  number of words of $c$ whose ball contains $u$ was computed from Hamming distances alone. It is
  $6$ for $u\in c$, $1$ for $u$ in a block of $N(c)$ and $0$ otherwise, as in the lemma.
\item \emph{Distributions.} A separate, vectorised enumeration of all $76\,904\,685$ ways to place
  $32$ words in nine blocks ($95\,548\,245$ for $33$), using the coefficients computed above,
  yields exactly the same $306$ ($1\,194$) distributions.
\item \emph{Symmetry reduction.} Each of these distributions is equivalent to one of the
  representatives, and each representative occurs.
\item \emph{The exact model.} A $36$-word covering code found independently satisfies the model
  for its own distribution, and the same code with one word removed does not.
\item \emph{Known values.} The same pipeline reproduces $K(2,3,1)=12$ and $K(2,4,1)=20$: it
  reports no code with $11$ and $19$ words, and finds codes with $12$ and $20$ words, which an
  independent brute-force checker confirms.
\end{itemize}

\begin{remark}
All three solvers use floating-point arithmetic, so each verdict is evidence rather than an
exact certificate; they come from different teams and agree wherever they finished.
\end{remark}

\begin{remark}
For $M=34$ there are $4\,266$ distributions, $93$ up to symmetry. SCIP finds $84$ of these
infeasible; $9$ were still open after $300$ seconds each. Settling them would give
$K(2,5,1)\ge 35$.
\end{remark}

\paragraph{Code and data.} All programs and logs are in the public repository
\begin{center}\url{https://github.com/Noah-Parsons/covering-codes}\end{center}
\begin{itemize}
\item the search: \texttt{src/split\_t2.py};
\item the Gurobi and HiGHS rechecks: \texttt{src/recheck\_t2\_gurobi.py},
  \texttt{src/recheck\_t2\_highs.py};
\item the checks of the reduction: \texttt{src/verify\_t2\_distributions.py};
\item logs: \texttt{results/}.
\end{itemize}

\begin{thebibliography}{9}

\bibitem{bow} R.~Bertolo, P.~R.~J. Östergård and W.~D. Weakley, An updated table of
binary/ternary mixed covering codes, \emph{J. Combin. Des.} 12 (2004), 157--176.

\bibitem{cohen-book} G.~Cohen, I.~Honkala, S.~Litsyn and A.~Lobstein, \emph{Covering Codes},
North-Holland, Amsterdam, 1997.

\bibitem{hr} H.~Hämäläinen and S.~Rankinen, Upper bounds for football pool problems and mixed
covering codes, \emph{J. Combin. Theory Ser.~A} 56 (1991), 84--95.

\bibitem{keri} G.~Kéri, Tables for bounds on covering codes,
\url{https://old.sztaki.hu/~keri/codes/index.htm} (accessed 16 October 2026).

\bibitem{lqs} W.~Lang, J.~Quistorff and E.~Schneider, New results on integer programming for
codes, \emph{Congr. Numer.} 188 (2007), 97--107.

\bibitem{oh} P.~R.~J. Östergård and H.~O. Hämäläinen, A new table of binary/ternary mixed
covering codes, \emph{Des. Codes Cryptogr.} 11 (1997), 151--178.

\end{thebibliography}
\end{document}
